\subsection{Anzahl der Lösungen erkennen}
\label{subsec:linear_systems_solution_criteria}

Sei $A$ die Koeffizientenmatrix eines LGS mit $n$ Unbekannten und
$[A\mid b]$ seine erweiterte Matrix. Nach der Vorwärtselimination wird zuerst
geprüft, ob eine Zeile der Form
\[
  \left[\begin{array}{cccc|c}
    0&0&\cdots&0&c
  \end{array}\right],
  \qquad c\neq0,
\]
auftritt. Sie entspricht der unmöglichen Gleichung $0=c$.

\begin{TheoremBox}{Lösbarkeits- und Rangkriterium}{linear_systems_rank_criterion}
  Für ein LGS $Ax=b$ mit $n$ Unbekannten gilt:
  \[
    \begin{array}{rcll}
      \operatorname{rang}(A)
      &<&\operatorname{rang}([A\mid b])
      &\Longleftrightarrow\quad \text{keine Lösung},\\[0.7em]
      \operatorname{rang}(A)
      &=&\operatorname{rang}([A\mid b])=n
      &\Longleftrightarrow\quad \text{genau eine Lösung},\\[0.7em]
      \operatorname{rang}(A)
      &=&\operatorname{rang}([A\mid b])<n
      &\Longleftrightarrow\quad \text{unendlich viele Lösungen}.
    \end{array}
  \]
\end{TheoremBox}

Der Rangvergleich prüft zunächst die Verträglichkeit der Gleichungen. Sind die
Ränge gleich, entscheidet die Zahl der Pivotpositionen relativ zur Zahl $n$
der Unbekannten, ob freie Variablen vorhanden sind.

\begin{center}
  \begin{tikzpicture}[
  >=Latex,
  box/.style={
    draw=black!25,fill=black!2,rounded corners=2pt,
    text width=4.25cm,minimum height=1.2cm,align=center,font=\small
  },
  result/.style={
    draw=THMBlue!55,fill=THMBlue!4,rounded corners=2pt,
    text width=4.35cm,minimum height=1.15cm,align=center,font=\small\bfseries
  },
  arrow/.style={->,line width=0.8pt,THMBlue}
]
  \node[box] (contradiction)
    {Widerspruchszeile?\\$[,0\;\cdots\;0\mid c,]$, $c\neq0$};
  \node[result,right=1.2cm of contradiction] (none)
    {keine Lösung\\$\operatorname{rang}(A)<\operatorname{rang}([A\mid b])$};
  \node[box,below=0.9cm of contradiction] (free)
    {Mindestens eine freie Variable?\\$\operatorname{rang}(A)<n$};
  \node[result,right=1.2cm of free] (many)
    {unendlich~viele\\Lösungen\\$n-\operatorname{rang}(A)$ Parameter};
  \node[result,below=0.9cm of free] (one)
    {genau eine Lösung\\$\operatorname{rang}(A)=n$};

  \draw[arrow] (contradiction)--node[above,font=\scriptsize]{ja}(none);
  \draw[arrow] (contradiction)--node[left,font=\scriptsize]{nein}(free);
  \draw[arrow] (free)--node[above,font=\scriptsize]{ja}(many);
  \draw[arrow] (free)--node[left,font=\scriptsize]{nein}(one);
\end{tikzpicture}

\end{center}

\begin{ExampleBox}{Drei letzte Zeilen, drei Aussagen}{linear_systems_three_last_rows}
  Für ein LGS mit drei Unbekannten können nach der Elimination beispielsweise
  folgende Situationen auftreten:
  \[
    \begin{array}{c@{\qquad}l}
      [,0\;0\;0\mid -4,]
        &\Rightarrow 0=-4:\quad \text{keine Lösung},\\[0.7em]
      [,0\;0\;3\mid -63,]
        &\Rightarrow x_3=-21:\quad \text{Pivot in jeder Variablenspalte},\\[0.7em]
      [,0\;0\;0\mid 0,]
        &\Rightarrow 0=0:\quad \text{keine neue Bedingung}.
    \end{array}
  \]
  Im letzten Fall entscheidet die Zahl der vorherigen Pivotpositionen, ob eine
  oder mehrere Variablen frei sind.
\end{ExampleBox}

\begin{RemarkBox}
  Die Zahl der Nichtnullzeilen der \emph{erweiterten} Matrix allein ist kein
  Kriterium für unendlich viele Lösungen: Eine Widerspruchszeile ist dort eine
  Nichtnullzeile. Deshalb wird stets zuerst die Lösbarkeit geprüft.
\end{RemarkBox}
